% Deck figure n12: Clifford gates are not applied to the data (README Section 3.5, "Clifford frame").
%   left    one step of the kicked Ising circuit as a sketch on six qubits in a line: an X rotation on every
%           qubit, then the layers of ZZ(pi/2) gates, which are Clifford gates
%   right   what is applied to the terms: the ZZ gates are absorbed into the generators of the rotations that
%           follow. On the line of the sketch X becomes Z X Z (Y Z at the two ends); signs are left out.
%   bottom  the 127-qubit kicked Ising circuit with 20 steps: 5,420 gates, of which 2,880 are Clifford ZZ gates
%           (5,420 - 2,540); 2,541 operations are applied to the terms: the 2,540 rotations and one operation
%           that only carries a truncation (README Section 3.3: 2,541 operations; Section 3.5 and the row
%           "Clifford frame" of the table in Section 4.7: 5,420 gates become 2,540 rotations).
\begin{tikzpicture}[x=1cm, y=1cm, font=\rmfamily\footnotesize]
  \tikzset{g1/.style={draw=black!70, line width=0.35pt, fill=genc, inner sep=0pt, minimum width=0.38cm, minimum height=0.38cm,
                      font=\rmfamily\scriptsize},
           zz/.style={draw=black!70, line width=0.35pt, fill=black!8, inner sep=0pt, minimum width=0.38cm, minimum height=0.92cm,
                      postaction={pattern=north east lines, pattern color=black!55}},
           rc/.style={draw=black!70, line width=0.35pt, anchor=north west, inner sep=0pt, minimum height=0.46cm,
                      text height=0.28cm, text depth=0.09cm, font=\rmfamily\footnotesize}}
  % ================= the circuit =================
  \node[frame, anchor=north west, minimum width=5.3cm, minimum height=5.4cm] at (0,0) {};
  \node[ttl, anchor=north] at (2.65,-0.06) {Kicked Ising Circuit, One Step};
  \foreach \k in {0,...,5}
  {
    \draw[line width=0.5pt, black!80] (0.35,{-1.15-0.54*\k}) -- (4.95,{-1.15-0.54*\k});
    \node[g1] at (1.05,{-1.15-0.54*\k}) {X};
  }
  \foreach \x/\k in {2.2/0, 2.2/3, 3.1/1, 3.1/4, 4.0/2}
    \node[zz] at (\x,{-1.42-0.54*\k}) {};
  \node[g1] at (0.6,-4.74) {};
  \node[inote, anchor=base west] at (0.82,-4.82) {X Rotations};
  \node[zz, minimum height=0.38cm] at (2.75,-4.74) {};
  \node[inote, anchor=base west] at (2.97,-4.82) {Clifford ZZ Layers};
  % ================= what is applied to the terms =================
  \node[frame, anchor=north west, minimum width=5.3cm, minimum height=5.4cm] at (8.6,0) {};
  \node[ttl, anchor=north] at (11.25,-0.06) {What Is Applied to the Terms};
  \foreach \k in {0,...,5}
    \draw[line width=0.5pt, black!80] (8.95,{-1.15-0.54*\k}) -- (13.55,{-1.15-0.54*\k});
  % a rotation box spans the qubits of its generator: \a top wire, \n number of wires
  \foreach \x/\a/\n in {9.9/0/2, 9.9/2/3, 11.25/0/3, 11.25/3/3, 12.6/1/3, 12.6/4/2}
    \node[g1, minimum height={0.54*(\n-1)*1cm+0.38cm}] at (\x,{-1.15-0.54*\a-0.27*(\n-1)}) {};
  \foreach \x/\k/\t in {9.9/0/Y, 9.9/1/Z, 9.9/2/Z, 9.9/3/X, 9.9/4/Z, 11.25/0/Z, 11.25/1/X, 11.25/2/Z,
                        11.25/3/Z, 11.25/4/X, 11.25/5/Z, 12.6/1/Z, 12.6/2/X, 12.6/3/Z, 12.6/4/Z, 12.6/5/Y}
    \node[font=\rmfamily\scriptsize, inner sep=0pt] at (\x,{-1.15-0.54*\k}) {\t};
  \node[g1] at (9.25,-4.74) {};
  \node[inote, anchor=base west] at (9.47,-4.82) {Rotations With Transformed Generators};
  % ================= between the two =================
  \draw[fat] (5.45,-2.5) -- (8.45,-2.5);
  \node[note] at (6.95,-1.3) {Absorbed Into the};
  \node[note] at (6.95,-1.63) {Generators of the};
  \node[note] at (6.95,-1.96) {Following Rotations};
  % ================= measured =================
  \node[rc, fill=black!10, minimum width=5.3cm] at (0,-5.65) {Gates in the Circuit};
  \node[rc, fill=black!10, minimum width=3.3cm] at (5.3,-5.65) {Clifford Gates Absorbed};
  \node[rc, fill=black!10, minimum width=5.3cm] at (8.6,-5.65) {Operations Applied to the Terms};
  \node[rc, fill=white, minimum width=5.3cm] at (0,-6.11) {5,420};
  \node[rc, fill=white, minimum width=3.3cm] at (5.3,-6.11) {2,880};
  \node[rc, fill=white, minimum width=5.3cm, text=dred, font=\rmfamily\bfseries\footnotesize] at (8.6,-6.11) {2,541};
  \node[note] at (6.95,-6.87) {2,541 = 2,540 Rotations and One Operation That Only Carries a Truncation};
\end{tikzpicture}
